\section{Symmetry accumulated between two visits to a coordinate}
\label{shear:sec:delayed-symmetry}

Intervening coordinate switches produce a symmetry that can be used at
an earlier conditioning time. Between two visits to a coordinate, they
make the available set and its conditional mass vector invariant in law
under every translation of one coordinate half. This yields an energy-loss
bound conditional on the earlier visit: the two half counts are already
fixed at that time, while the intervening switches supply the symmetry.

\begin{lemma}[Delayed half-space symmetry]
\label{shear:lem:delayed-symmetry}
Let $V=\mathbb F_2^d$, $N=2^d$, and fix a coordinate $i$. Let
$j_1,\ldots,j_{d-1}$ list all coordinates other than $i$, in any order.
Condition on a sigma field $\mathcal A$ which specifies a nonempty set
$B_0\subseteq V$ of size $h$ and a probability vector $p_0$ supported on
$B_0$. Independently of $\mathcal A$, let $S_r$ be independent uniform
matching switches in direction $e_{j_r}$, for $1\le r\le d-1$.
Define $B_r=S_rB_{r-1}$. On each matching edge, obtain $p_r$ by averaging
the two masses when both endpoints belong to $B_{r-1}$, and by moving
the mass with its unique available endpoint when exactly one endpoint
belongs to $B_{r-1}$; it is zero on edges with no available endpoint.
Put
\[
 a_r(x)=p_r(x)-h^{-1}\mathbf1_{B_r}(x),\qquad
 E_r=\sum_{x\in V}a_r(x)^2,
 \qquad H_b=\{x\in V:x_i=b\}\quad(b=0,1).
\]
For each $z\in V$ with $z_i=0$, let $T_zx=x+x_i z$. Conditional on
$\mathcal A$, the terminal pair $(B_{d-1},a_{d-1})$ has the same law as
\begin{equation}
 \bigl(T_zB_{d-1},\ a_{d-1}\circ T_z^{-1}\bigr).
 \label{shear:eq:delayed-symmetry-invariance}
\end{equation}
The counts $h_b=|B_r\cap H_b|$ are independent of $r$ and are
$\mathcal A$-measurable. If a further independent fair matching switch
in direction $e_i$ produces the energy $E_d$ by the same rule, then
\begin{equation}
 \mathbb E[E_{d-1}-E_d\mid\mathcal A]
 \ge\frac{\min(h_0,h_1)}N\,
             \mathbb E[E_{d-1}\mid\mathcal A].
 \label{shear:eq:delayed-symmetry-loss}
\end{equation}
The statement includes $d=1$, when the intermediate schedule is empty.
\end{lemma}

\begin{proof}
Fix $z$ with $z_i=0$, and set
\[
 z_0=0,\qquad z_r=\sum_{s=1}^r z_{j_s}e_{j_s}
       \quad(1\le r\le d-1).
\]
Thus $z_{d-1}=z$. Replace the $r$th switch by
\[
 S'_r=T_{z_r}S_rT_{z_{r-1}}^{-1}.
\]
We verify that the replaced switches have exactly the original joint law.
For $j\ne i$, $T_w e_j=e_j$, so conjugation by $T_w$ maps the matching
subgroup in direction $e_j$ onto itself. Also
\[
 T_{z_r}T_{z_{r-1}}^{-1}(x)
   =x+x_i z_{j_r}e_{j_r}
\]
is a fixed member of that matching subgroup: it either swaps or fixes
each pair, since $x_i$ is constant on a coordinate-$j_r$ pair. Therefore
$S_r\mapsto S'_r$ is a bijection of this subgroup. It depends only on
$r$ and $z$, so the $S'_r$ remain independent and uniform conditional
on $\mathcal A$.

Run the available-set and mass recursion with these replaced switches,
starting from the same $(B_0,p_0)$. We claim that at stage $r$ its state is
$(T_{z_r}B_r,p_r\circ T_{z_r}^{-1})$. This holds initially because
$T_{z_0}$ is the identity. Conjugation carries the matching edges at the
previous stage bijectively to matching edges. An edge with two available
endpoints is averaged in either recursion, and a subsequent fixed swap
does not change its equal output masses. An edge with one available
endpoint carries its mass to the image of that endpoint under the
replaced switch. The assertion is immediate on edges with no available
endpoint. These three cases prove the induction. Centering by the
uniform mass on the available set transforms in the same way. Since
the two switch arrays have the same conditional law, the terminal
invariance follows.

Every intermediate switch preserves $x_i$, proving the assertion about
the counts. The squared-energy decrease in the final coordinate-$i$
layer is
\[
 L(B,a)=\frac12
 \sum_{\substack{\{x,y\}\subseteq B\\x+y=e_i}}
                 (a(x)-a(y))^2,
 \qquad (B,a)=(B_{d-1},a_{d-1}).
\]
Indeed, transport preserves squared norm and averaging two entries
decreases it by half the square of their difference. We may average
this expression after applying an independent uniform $T_z$, by the
conditional invariance already proved. Each pair
$x\in B\cap H_0$, $y\in B\cap H_1$ becomes a coordinate-$i$ pair
for exactly one of the $N/2$ possible values of $z$. Consequently
\[
 \mathbb E_z L(T_zB,a\circ T_z^{-1})
 =\frac1N\sum_{x\in B\cap H_0,\ y\in B\cap H_1}
                         (a(x)-a(y))^2.
\]
The entries of $a$ sum to zero. With
$A=\sum_{x\in B\cap H_0}a(x)$, the last double sum equals
\[
 h_1\sum_{B\cap H_0}a(x)^2+
 h_0\sum_{B\cap H_1}a(y)^2+2A^2
 \ge\min(h_0,h_1)E_{d-1}.
\]
Taking conditional expectations proves
\eqref{shear:eq:delayed-symmetry-loss}. Empty halves give a zero lower
bound and require no separate division. When $d=1$, the same calculation
has only $z=0$ and remains valid.
\end{proof}

In a hidden-card application, $B_0$ and $p_0$ are the available set and
its conditional mass vector at the conditioning time. They are functions
of the observed path prefix. Conditioning additionally on all past
switches fixes these two functions; it does not redefine $p_0$ as a
conditional law given every card position. Future matching coins are
still independent and fair, so the lemma applies to the stated
transport-and-averaging recursion. If the coordinate schedule is any
periodic permutation of the $d$ coordinates, take $\mathcal A$ immediately
after the previous visit to $i$ and apply the lemma to the intervening $d-1$
layers. Reversing the coordinate order is included in the same statement.
